%!TEX root = ../optim-ml/OptimML-PAC.tex



Assume the observations $\Dd \eqdef (x_i,y_i)_{i=1}^n$ are independent and identically distributed copies of a random pair $(X,Y)$ taking values in $\Xx\times\Yy$. From this sample, we seek a predictor $f : \Xx \rightarrow \Yy$ with risk close to the smallest achievable value:
\eq{
	L(f) \eqdef \EE( \ell(Y,f(X)) ) = \int_{\Xx \times \Yy} \ell(y,f(x)) \d\PP_{X,Y}(x,y)
}
where $\ell : \Yy \times \Yy \rightarrow \RR$ is some loss function. 
%
To estimate this minimizer, we select a class of functions $\Ff$ and minimize the empirical risk
\eq{
	\hat f \in \uargmin{f \in \Ff} \hat L(f) \eqdef \frac{1}{n} \sum_{i=1}^n \ell(y_i,f(x_i)).
}	
We assume a measurable minimizer exists. It is random because it depends on the sample $\Dd$.

\begin{exmp}[Regression and classification]
	In regression, $\Yy=\RR$ and the most common choice of loss is $\ell(y,z) = (y-z)^2$.
	For binary classification, $\Yy=\{-1,+1\}$ and the ideal choice is the 0-1 loss $\ell(y,z)=1_{y \neq z}$.
	%	
	For many useful classes $\Ff$, minimizing the empirical 0--1 risk is computationally difficult. A common alternative uses a surrogate loss of the form $\ell(y,z) = \Gamma(-yz)$ where $\Gamma$ is a convex function upper-bounding $1_{\RR^+}$, which makes empirical risk minimization convex when the real-valued score class is convex.
	\end{exmp}
	
PAC learning seeks a bound on the excess risk $L(\hat f)-\inf(L)\geq0$ that holds with probability at least $1-\delta$. The bound depends on the sample size $n$ and confidence parameter $\delta$. For the excess risk to vanish, the approximation class must become sufficiently rich while its statistical complexity remains controlled. Quantitative approximation rates require distributional assumptions, although some consistency results hold without them.


%%%%%%%%%%%%%%%%%%%%%%%%%%%%%%%%%%%%%%%%%%%%%%%%%%
\subsection{Nonparametric Models and Calibration}

If $\Ff$ contains all measurable functions, a minimizer $f^\star$ of $L$ is called a Bayes estimator: it achieves the smallest possible population risk.


%%%
\paragraph{Risk decomposition}

Denoting 
\eq{
	\al(z|x) \eqdef \EE_Y(\ell(Y,z)|X=x) = \int_\Yy\ell(y,z)\d\PP_{Y|X}(y|x)
} 
the average loss of predicting $z\in\Yy$ at input $x\in\Xx$, the risk decomposes as
\eq{
	L(f) = \int_\Xx \left[ \int_\Yy \ell(y,f(x)) \d \PP_{Y|X}(y|x) \right] \d \PP_X(x) = \int_\Xx \al(f(x)|x) \d \PP_X(x)
}
Consequently, $f^\star$ can be chosen separately at each input $x$ by solving
\eq{
	f^\star(x) = \uargmin{z} \al(z|x).
}

\begin{exmp}[Regression]
For squared-loss regression with $\EE(Y^2)<\infty$, one has $f^\star(x)=\EE(Y|X=x) = \int_\Yy y \d \PP_{Y|X}(y|x)$, which is the conditional expectation.
\end{exmp}

\begin{exmp}[Classification]
For classification applications, where $\Yy = \{-1,1\}$, it is convenient to introduce $\eta(x) \eqdef \PP_{Y|X}(y=1|x) \in [0,1]$.
%
If the two classes are separable, then $\eta(x) \in \{0,1\}$ for $\PP_X$-almost every $x$.
%
For the 0-1 loss $\ell(y,z)=1_{y\neq z}=1_{(0,+\infty)}(-yz)$, one may choose $f^\star=\sign(2\eta-1)$, with either class selected at a tie and $L(f^\star) = \EE_X( \min(\eta(X),1-\eta(X)) )$.
%
The conditional probability $\eta$ is unknown and must be estimated from $\Dd$. Direct 0--1 optimization is computationally difficult for many useful classes $\Ff$.
%
For a margin loss $\ell(y,z)=\Ga(-yz)$, an unconstrained Bayes score satisfies
\eq{
	f^\star(x) = \uargmin{z} \Big( \al(z|x) = \eta(x) \Gamma(-z) + (1-\eta(x)) \Gamma(z) \Big), 
}
Thus the optimal score is a function of $\eta(x)$. This argmin notation presumes attainment; for logistic loss and $\eta\in\{0,1\}$, the infimum is approached only as the score tends to the appropriate infinity.
\end{exmp}

%%%
\paragraph{Calibration in the classification setup}

We would like the classifier $\sign(f^\star)$ to agree with a Bayes classifier for the 0--1 loss, namely $\sign(2\eta-1)$ away from ties. This property is called classification calibration. It does not require the scores $f^\star$ and $2\eta-1$ themselves to agree. The following result characterizes calibration for finite convex margin losses.

\begin{prop}
	A loss  $\ell$ associated to a finite, nonnegative convex $\Gamma$ is calibrated if and only if $\Gamma$ is differentiable at $0$ and $\Gamma'(0)>0$.
\end{prop}

Both hinge and logistic loss satisfy this calibration criterion.
% 
Let $L_\Gamma$ be the risk for $\ell(y,z)=\Gamma(-yz)$ and let $L_{0\text{--}1}$ be classification risk, with a fixed tie-breaking convention. Quantitative calibration bounds have the form
\eql{\label{eq-calibration-control}
	0 \leq L_{0\text{--}1}(\sign f)-\inf L_{0\text{--}1} \leq \Psi( L_{\Gamma}(f) - \inf L_{\Gamma} )
}
for some increasing function $\Psi : \RR^+ \rightarrow \RR^+$ satisfying $\Psi(0)=0$. Such a control ensures in particular that if $f^\star$ minimizes $L_{\Gamma}$, its induced classifier has Bayes-optimal 0--1 risk. At $\eta=1/2$, either class is optimal; away from ties, its sign agrees almost surely with that of $2\eta-1$.
%
For hinge loss, the calibration bound holds with $\Psi(r)=r$. For logistic loss, it holds with $\Psi(s)=\sqrt{2s}$.


%%%%%%%%%%%%%%%%%%%%%%%%%%%%%%%%%%%%%%%%%%%%%%%%%%
%%%%%%%%%%%%%%%%%%%%%%%%%%%%%%%%%%%%%%%%%%%%%%%%%%
\subsection{PAC bounds}

%%%
\paragraph{Approximation--estimation decomposition.}

For a class $\Ff$ of functions, the excess risk of the empirical estimator
\eq{
	\hat f \in \uargmin{f \in \Ff} \hat L(f)
%	\qandq
%	f^\star \eqdef \uargmin{f \text{ measurable}} L(f)	
} 
decomposes into a random estimation error and a deterministic approximation error:
\eql{\label{eq-bias-var-split}
	L(\hat f) - \inf L = 
		\Big[ L(\hat f)  - \inf_{\Ff} L \Big]
		+
		\Aa(\Ff)
		\qwhereq 
		\Aa(\Ff) \eqdef \Big[ \inf_{\Ff} L  - \inf L \Big].
}
This separates the effects of sampling and model approximation.

Enlarging a nested class $\Ff$ reduces approximation error but typically increases the estimation-error bound. Choose the class size as a function of $n$ to balance these effects and avoid overfitting.

%%%%%%%%%%%%%%%%%%%%%%%%%%%%%%%%%%%%%%%%%%%%%%%%%%
\paragraph{Approximation error.}

A quantitative approximation bound requires assumptions on $f^\star$. No-free-lunch results rule out a uniform distribution-free approximation rate over all measurable targets.
%
The following example illustrates such a bound.

%%%%
\begin{exmp}[Linearly parameterized functionals]
A popular class of functions consists of linearly parameterized maps of the form
\eq{
	f(x) = f_w(x) = \dotp{\phi(x)}{w}_\Hh
}
where $\phi : \Xx \rightarrow \Hh$ maps inputs to a Hilbert feature space. For $\Xx=\RR^p$, the identity map $\phi(x)=x$ recovers ordinary linear models. One can also consider for instance polynomial features, $\phi(x)=(1, x_1,\ldots,x_p,x_1^2,x_1 x_2, \ldots)$, giving rise to polynomial regression and polynomial classification boundaries.
%
Restrict the parameters to the ball $\Ff = \enscond{f_w}{\norm{w}_\Hh \leq R}$. Suppose $f^\star=f_{w^\star}$ for some $w^\star\in\Hh$, the loss $\ell(y,\cdot)$ is uniformly $Q$-Lipschitz, and $\EE\norm{\phi(X)}_\Hh<\infty$. Projecting $w^\star$ onto the radius-$R$ ball gives
\eq{
	\Aa(\Ff)  \leq Q \EE(\norm{\phi(X)}_\Hh) \max( \norm{w^\star}_\Hh-R, 0 ).
}
\end{exmp}

%%%%
\begin{rem}[Connections with RKHS]
Note that this lifting actually corresponds to using functions $f$ in a reproducing kernel Hilbert space, denoting
\eq{
	\norm{f}_k \eqdef \uinf{w \in \Hh} \enscond{ \norm{w}_\Hh }{ f = f_w }
}
with associated kernel $k(x,x') \eqdef \dotp{\phi(x)}{\phi(x')}_\Hh$. We use the feature-map representation directly below.
\end{rem}


%%%%%%%%%%%%%%%%%%%%%%%%%%%%%%%%%%%%%%%%%%%%%%%%%%
\paragraph{Estimation error.}

Concentration inequalities control sampling fluctuations, while a bound on the complexity of $\Ff$ makes this control uniform over the model class.

The estimation error is controlled by uniform deviations between $L$ and $\hat L$. Let $g\in\Ff$ attain $\inf_{\Ff}L$, assuming for simplicity that such a minimizer exists. Then
\eql{\label{eq-estimation-split}
	L(\hat f) - \inf_{\Ff} L = 
	\Big[L(\hat f) - \hat L(\hat f)\Big] 
	+
	\Big[\hat L(\hat f) - \hat L(g)\Big] 
	+
	\Big[\hat L(g)-L(g)\Big]
	\leq 
	2 \sup_{f \in \Ff}|\hat L(f) - L(f)|
}
since $\hat L(\hat f) - \hat L(g) \leq 0$.
%
Define the two one-sided deviations
\eq{\De_+(\Dd)=\sup_{f\in\Ff}(L(f)-\hat L(f)),\qquad
\De_-(\Dd)=\sup_{f\in\Ff}(\hat L(f)-L(f)).}
The preceding decomposition gives the sharper bound $L(\hat f)-\inf_\Ff L\leq\De_++\De_-$. Assume all suprema below are measurable (or use outer expectations).

\begin{prop}[McDiarmid concentration]\label{prop-mc-diarmid}
Assume $0\leq\ell(Y,f(X))\leq\ell_\infty$ almost surely, uniformly over $f\in\Ff$. With probability at least $1-\delta$, both signs satisfy
\eql{\label{eq-mc-diarmid}
\De_\pm(\Dd)-\EE_\Dd\De_\pm(\Dd)
\leq\ell_\infty\sqrt{\frac{\log(2/\delta)}{2n}}.
}
\end{prop}

To bound $\EE_\Dd\De_\pm(\Dd)$, we need to control the complexity of $\Ff$. VC dimension gives one approach, but for norm-constrained linear models it can yield loose bounds that depend on the ambient dimension. Rademacher complexity gives a finer measure for a class $\Gg$ of functions from $\Xx \times \Yy$ to $\RR$:
\eq{
	\Rr_n(\Gg) \eqdef \EE_{\epsilon,\Dd}\Big[
		\sup_{g \in \Gg} \frac{1}{n}\sum_{i=1}^n \epsilon_i g(x_i,y_i)
	\Big]
}
where $\epsilon_i$, independent of the data, are independent Rademacher random variables (i.e. $\PP(\epsilon_i = \pm 1) = 1/2$). The quantity $\Rr_n(\Gg)$ depends on the distribution of $(X,Y)$.

Apply this complexity measure to the loss class $\Gg = \ell[\Ff]$, defined by
\eq{
	\ell[\Ff] \eqdef \enscond{ (x,y) \in \Xx \times \Yy \mapsto \ell(y,f(x))  }{ f \in \Ff }, 
}
Symmetrization gives the following bounds.

\begin{prop}\label{prop-control-rademacher}
	One has
	\eq{
		\EE\De_\pm(\Dd)\leq2\Rr_n(\ell[\Ff]).
	}
	If $\ell$ is $Q$-Lipschitz with respect to its second variable, one furthermore has $\Rr_n(\ell[\Ff]) \leq Q \Rr_n(\Ff)$ (here the class of functions only depends on $x$).
\end{prop}

Combining~\eqref{eq-bias-var-split}, \eqref{eq-estimation-split}, and Propositions~\ref{prop-mc-diarmid} and \ref{prop-control-rademacher} gives the following result.

\begin{thm}
	Assume $\ell(y,\cdot)$ is uniformly $Q$-Lipschitz and $0\leq\ell\leq\ell_\infty$ uniformly over the class. Then, with probability at least $1-\de$
	\eql{\label{eq-final-pac}
		0 \leq L(\hat f) - \inf L \leq
			\ell_\infty\sqrt{\frac{2\log(2/\de)}{n}}
			+ 
			4 Q \Rr_n(\Ff)
			+ 
			\Aa(\Ff).		
	}
\end{thm}

\begin{exmp}[Linear models]
	In the case where $\Ff = \enscond{ \dotp{\phi(\cdot)}{w}_{\Hh} }{ \norm{w} \leq R}$
	where $\norm{\cdot}$ is some norm on $\Hh$, one has
	\eq{
		\Rr_n(\Ff)=\frac{R}{n}\EE_{\epsilon,\Dd}\norm{\sum_i\epsilon_i\phi(x_i)}_*
	}
	where $\norm{\cdot}_*$ is the dual norm, defined by
	\eq{
		\norm{u}_* \eqdef \usup{ \norm{w} \leq 1} \dotp{u}{w}_\Hh.
	}
	When $\norm{\cdot}=\norm{\cdot}_\Hh$ is the Hilbert norm, the expression simplifies further:
	$\norm{u}_* = \norm{u}_\Hh$ and
	\eq{
		\Rr_n(\Ff) \leq \frac{R \sqrt{\EE(\norm{\phi(X)}_\Hh^2)} }{\sqrt{n}}.
	}
	The bound has no explicit dependence on feature dimension and also applies to infinite-dimensional RKHS feature maps with finite second moment.
	%
	For fixed $R$ and fixed loss constants, the estimation bound in~\eqref{eq-final-pac} decays as $1/\sqrt{n}$; the approximation error remains a separate term.
	%
	The radius $R$ controls the balance between approximation and estimation. In practice, it can be selected by cross-validation.
\end{exmp}

\begin{exmp}[Application to SVM classification]
The 0--1 loss is not Lipschitz in a real-valued score. For a margin $\rho>0$, use the bounded ramp loss
\eq{\Gamma_\rho(s)=\min\{1,\max\{0,1+s/\rho\}\}.}
It is $1/\rho$-Lipschitz and satisfies
\eq{1_{[0,+\infty)}(s)\leq\Gamma_\rho(s)\leq\max\{1+s/\rho,0\}.}
The left inequality counts score ties as errors, so it also bounds the risk of any fixed tie-breaking classifier. For $\Ff_R=\{f_w:\norm{w}_\Hh\leq R\}$, the one-sided concentration and symmetrization bounds imply, simultaneously for every $f\in\Ff_R$, with probability at least $1-\delta$,
\eq{
 L_{0\text{--}1}(\sign f)\leq
 \frac1n\sum_i\max\{1-y_if(x_i)/\rho,0\}
 +\frac{2R}{\rho}\sqrt{\frac{\EE\norm{\phi(X)}_\Hh^2}{n}}
 +\sqrt{\frac{\log(1/\delta)}{2n}}.
}
This bound applies to a data-dependent predictor, including an empirical hinge-loss minimizer in the ball. It does not require minimizing the nonconvex ramp loss. In practice, the norm constraint is often replaced by a penalty:
\eql{\label{eq-svm-optim}
\umin{w}\frac1n\sum_i\max\{1-y_if_w(x_i),0\}+\lambda\norm{w}_\Hh^2,\qquad\lambda>0.
}
This is the kernel support vector machine.
\end{exmp}

\begin{rem}[Resolution using the kernel trick]
	The kernel trick reduces problems such as~\eqref{eq-svm-optim}, of the form
	\eql{\label{eq-kernel-method}
		\umin{w \in \Hh} \frac{1}{n} \sum_{i=1}^n \ell(y_i, \dotp{w}{\phi(x_i)} ) + \lambda \norm{w}_\Hh^2
	}
	to optimization over $n$ coefficients, even when the feature space is infinite dimensional. The representer theorem gives
	$w^\star = \sum_{i=1}^n c_i^\star \phi(x_i)$ for some $c^\star \in \RR^n$: an orthogonal component would leave the predictions unchanged and strictly increase the positive norm penalty.
	% 
	Substituting this representation into~\eqref{eq-kernel-method} reduces the data dependence to the empirical kernel matrix $K = ( k(x_i,x_j)\eqdef \dotp{\phi(x_i)}{\phi(x_j)} )_{i,j=1}^n \in \RR^{n \times n}$ since it reads
	\eq{
		\umin{c \in \RR^n} \frac{1}{n} \sum_{i=1}^n \ell(y_i, (K c)_i) + \lambda \dotp{K c}{c}
	}
	From any optimal $c^\star$ of this convex program (which need not be strictly convex when $K$ is singular), the estimator is evaluated as
	$f_{w^\star}(x)=\dotp{w^\star}{\phi(x)} = \sum_i c_i^\star k(x_i,x)$. 
\end{rem} 